% Generated draft: requires mathematical and source review.
\documentclass[11pt]{article}
\usepackage[margin=25mm]{geometry}
\usepackage{fontspec}
\setmainfont{Latin Modern Roman}
\usepackage{amsmath,amssymb,mathtools,unicode-math}
\setmathfont{Latin Modern Math}
\usepackage{enumitem,needspace,xurl,hyperref}
\setlength{\parindent}{0pt}
\setlength{\parskip}{5pt}
\setlength{\emergencystretch}{3em}
\newcommand{\N}{\mathbb{N}}
\newcommand{\Z}{\mathbb{Z}}
\newcommand{\Q}{\mathbb{Q}}
\newcommand{\R}{\mathbb{R}}
\newcommand{\C}{\mathbb{C}}
\newcommand{\F}{\mathbb{F}}
\newcommand{\A}{\mathbb{A}}
\newcommand{\PP}{\mathbb{P}}
\newcommand{\E}{\mathbb{E}}
\newcommand{\eps}{\varepsilon}
\newcommand{\dd}{\,\mathrm{d}}
\begin{document}
\section*{UnsolvedMath ID: 801097}
\textbf{Source problem code:} catalogue entry
\subsection{Definitions and mathematical statement}
For \(c,\delta\in\Q\), define
\[
f_{c,\delta}\colon\Q^2\longrightarrow\Q^2,\qquad f_{c,\delta}(x,y)=(y,y^2+c-\delta x).
\]
A point \(P\in\Q^2\) is periodic if \(f_{c,\delta}^{\circ n}(P)=P\) for some integer \(n\geq1\); its exact period is the least such integer.

The supplied conjecture asserts that there exists a finite set \(E\subset\Q\) such that, for every \(\delta\in\Q\setminus E\), every \(c\in\Q\), and every periodic point \(P\in\Q^2\), the exact period of \(P\) divides \(2\). Equivalently,
\[
\exists E\subset\Q\text{ finite}\quad
\forall\delta\in\Q\setminus E\quad
\forall c\in\Q\quad
\forall P\in\Q^2,\qquad
\bigl(\exists n\geq1:\ f_{c,\delta}^{\circ n}(P)=P\bigr)
\Longrightarrow f_{c,\delta}^{\circ2}(P)=P.
\]
Here \(n\) ranges over integers, and the exceptional set \(E\) is independent of \(c\). The supplied statement imposes no restriction \(\delta\neq0\), so the parameter domain is retained exactly as given.
\subsection{Short English statement}
The conjecture asserts that, outside one finite set of rational values of \(\delta\), every rational periodic point of every map in this family, for every rational \(c\), has exact period \(1\) or \(2\).
\subsection{Sources}
\begin{itemize}
\item Local UnsolvedMath registry: permanent ID \(801097\), source problem code \texttt{catalogue entry}, supplied title \emph{Henon Maps: Open Problems --- Conjecture 49}. The registry labels the problem open; this is a supplied status, not an independently verified assessment.
\item Supplied attribution: Berger, Bedford, \emph{Henon Maps: A List of Open Problems} (2023), Conjecture 49. Supplied source URL: \url{https://arxiv.org/abs/2312.03907}. The provenance records access on 2026-10-16 and extraction by \texttt{arxiv\_latex}; these are supplied provenance details, not a claim of independent source consultation.
\end{itemize}
\end{document}
